% MT052 / explanation — canonical editable content.
生成式分类器需要知道两件事：每类内部的数据是什么样，以及每类有多常见。以前者\(p(x\mid c)\)描述新输入像不像某类，以后者\(P(c)\)表示先验，再比较两者乘积。



\begin{itemize}

\item \textbf{A正确。} 正是在每类内部估计特征分布。例如Gaussian NB学习该类各特征的均值和方差。

\item \textbf{B错误。} 类别先验是加入既有知识的直接入口。若已知某类在目标人群中的比例，可以在有依据时指定或调整\(P(c)\)；对参数也可加先验。

\item \textbf{C错误。} 有三类时，为三类各建立分布及先验，再比较三个分数即可，没有只能二分类的限制。

\item \textbf{D正确。} 分布假设决定哪些输入在模型看来常见。若数据是二值出现信息，却用不合适的连续形状强行拟合，估计的概率和分类表现都可能受影响。

\item \textbf{E错误。} 它只比较\(p(x\mid c)\)，漏了\(P(c)\)。罕见类的密度略高，仍可能不足以抵消它很低的先验。只有先验相等等条件下，单独比较密度才等价。

\end{itemize}


类条件密度与后验的数值比较见\StudyRef{Answers}{mt040}{MT040}。


\textbf{出处：} 2023B期中 Q1，原卷第2页（5分）


\textbf{答案依据：} 2023B随卷解答 Q1，随卷答案AD。
